% SPDX-FileCopyrightText: Copyright (c) 2026 NVIDIA CORPORATION & AFFILIATES. All rights reserved.
% SPDX-License-Identifier: Apache-2.0
%
% ROLE: the learned algorithm, completely enough that a router engineer could reimplement it:
%   how it plugs into the router, its inputs and exclusions, the default cost it starts from, the
%   features, the utility and decision rule (with pseudocode), why a context term and which form,
%   set-size invariance, parity with the default router, the nested-logit option, the
%   sticky-session comparators, porting notes, and the model ladder (which arms exist and run).
% CLAIMS: (1) the headline arms read only router-observable inputs (no AIS, no output length);
%   (2) theta = -e0 reproduces the default router (unit tests; replay 88/88 with reservoir ties);
%   (3) a decision costs O(N d + r d) arithmetic plus an O(N log N) sort of the candidates;
%   (4) every v1/v2 feature except the identity-bound ones, and every context source, is invariant
%       under duplicating the candidate set.
% EVIDENCE: WT/lib/router-plugins/builtin/src/learned_choice/{FEATURES.md,features.rs,mod.rs,
%   parameters.rs,tests.rs}; WT/lib/router-plugins/builtin/src/{choice.rs,session_map.rs,
%   sticky_session.rs}; WT/lib/router-plugins/builtin/src/default/{scorer.rs,picker.rs};
%   WT/lib/kv-router/src/plugins/worker_selection.rs; CR/facts/integration.json parity;
%   CR/facts/build_rust.json (tests, smoke); CR/audits/build/leakage-features-r{0,1}.md;
%   CR/CONTRACT.md (learned-choice, sticky-session, A15, A16, A17); CR/literature/LESSONS.md LR-03,
%   LR-05, LR-06, LR-07, LR-13, LR-14, LR-15; CR/facts/pilot.json feature_scaling;
%   WT/notes/learned-routing/campaign/PLAN.md "Learned model ladder".
% STATUS: final. Feature set v1 frozen at calibration; v2 used by the headline arm; v3 lives on the
%   separate AIS branch (CR/facts/ais_features.json). Not measured: per-decision latency.
% LAYOUT (2026-10-05 restructure): implementation details (provider, parsing, fig:yaml, cost
%   details, leakage audits, porting notes, per-feature details, alg:decide, draw disciplines,
%   tab:sources, pooled vs named context, set-size and parity test details, sec:nested,
%   sec:sticky) live in sections/appendix-policy.tex (app:policy); the drift-check details in
%   sections/appendix-tuning.tex (sec:drift); the v3 features in sections/appendix-ais-league.tex.

\section{The learned policy}
\label{sec:model}

\lc{} takes over the final step of Dynamo's worker selection. Given the eligible candidates for
one request and the signals the router already tracks for them, it scores each candidate with a
linear utility over a handful of features and routes to the highest score. Three constraints
shape the design. It reads only inputs the live router has, so the coefficients tuned in replay
can run unchanged in a live router (\cref{sec:integration}). Its coefficients are shared by all
candidates and its features do not change when the candidate set is duplicated, so one coefficient
vector applies at every worker count, including counts that training never sees
(\cref{sec:set-size}). And one coefficient vector, $\theta = -\basis{0}$, reproduces the default
router, so tuning starts from the status quo and every learned coefficient reads as a deviation
from it (\cref{sec:parity}).\footnote{The normative specification, \file{FEATURES.md}, sits next
  to the policy's Rust code; this section follows it.}

\subsection{Router integration}
\label{sec:integration}

% src: WT/lib/kv-router/src/plugins/worker_selection.rs (WorkerScorer, WorkerPicker, factory per
%      partition); WT/lib/router-plugins/builtin/src/learned_choice/mod.rs (policy, provider, register)
In the API of \cref{sec:background}, \lc{} registers under the type string
\code{learned-choice} as one scorer and one picker. The scorer is the default router's own cost
function (\cref{eq:default-cost}) evaluated at default weights; its output becomes feature 0
(\cref{sec:features}). The picker (a \code{WorkerPicker}) computes the other features, the
utilities and the choice.
% src: learned_choice/mod.rs (required_worker_inputs = CACHE | LOAD); features.rs (accessors used)
It declares the plugin inputs \code{CACHE} and \code{LOAD} and reads exactly the inputs of
\cref{sec:decision}, the prefix hashes only for feature set \code{v2}.

% src: learned_choice/parameters.rs (deny_unknown_fields, validate); FEATURES.md "Parameters"
The coefficients are ordinary parameters in the policy YAML, so a new coefficient vector is a new
file rather than a rebuild. Each CMA-ES candidate in training is one such file, and the file that
training selects is the file a live router would load. \Cref{app:policy} shows such a file
(\cref{fig:yaml}) and gives the provider that instantiates the policy and its strict parsing rules.

\paragraph{Cost per decision.}
% src: learned_choice/mod.rs (decide: table.compute, effective_weights, dot per row);
%      choice.rs (Chooser::begin sorts rows by worker); FEATURES.md "Utility and decision" ("O(N·d + r·d)")
A decision costs $O(\Nw\dfeat + \rank\dfeat)$ arithmetic plus an $O(\Nw\log\Nw)$ sort; its
latency against \defaultcf{} was not measured (\cref{app:policy}).

\paragraph{Excluded inputs.}
% src: FEATURES.md "Excluded inputs"; CR/CONTRACT.md "learned-choice policy" ("Never used")
The policy never reads the request's expected output length, which offline replay fills with the
true output length, so reading it would give the policy future information that no live router
has. It never reads the modeled prefill backlog (the router's \code{PREFILL\_TIME} input) or any
other estimate from \ais{}, the performance model that times the simulated engines, so it stays
independent of the simulator it is trained in and needs nothing at runtime that a live router
lacks. Two independent audits checked these exclusions (\cref{app:policy}).
\Cref{sec:threats-scope} lists the host settings a deployment must match.
% Porting notes (exclusive-affinity override of theta6; features 0, 3, 4, 7 and the v2 key;
% seedless entropy) moved in full to app:policy; main survivor: discussion sec:threats-scope.

\subsection{Features}
\label{sec:features}

% src: FEATURES.md "Feature set v1"; features.rs (TOKEN_SCALE 8192, REQUEST_SCALE 32,
%      FALLBACK_KV_BLOCKS 65,536); CR/facts/setup.json (num_gpu_blocks 18,863)
\begin{table}[tbp]
  \centering\small
  \caption{Feature set \code{v1}, in order. $\Tch = 8192$ tokens, the engine's chunked-prefill
    batch; $\reqblocks_j = \max(\lceil \ISL_j/16 \rceil, 1)$. Feature 0 is computed by the default
    scorer itself at default weights, so $\theta = -\basis{0}$ reproduces the default router.}
  \label{tab:features}
  \begin{tabular}{@{}rll@{}}
    \toprule
    $f$ & Name & $x_{i,f}$ \\
    \midrule
    0 & \code{default\_logit\_scaled}   & $\cdef_i / \reqblocks_j$ at default weights \\
    1 & \code{overlap\_frac}            & $\ovl_i / \reqblocks_j$ \\
    2 & \code{new\_prefill\_tokens\_k}  & $\max(\ISL_j - \hit_i, 0) / \Tch$ \\
    3 & \code{active\_prefill\_tokens\_k} & $\Lpf_i / \Tch$ \\
    4 & \code{kv\_load\_frac}           & $\Ldec_i / \capb_i$ \\
    5 & \code{active\_requests\_s}      & $\Lreq_i / 32$ \\
    6 & \code{session\_affinity}        & 1 if $i$ served the previous request of $j$'s session, else 0 \\
    7 & \code{isl\_x\_prefill\_load}    & $(\ISL_j / \Tch)\, x_{i,3}$ \\
    \bottomrule
  \end{tabular}
\end{table}

\Cref{tab:features} defines the eight \code{v1} features, which have been frozen since
calibration began. They are built only from the inputs listed in \cref{sec:decision}.
\begin{itemize}
  \item \textbf{The anchor.} Feature 0 is the default cost divided by the request's block count.
    The division does not change the default's argmin, because $\reqblocks_j$ is the same for
    every candidate. Because the default scorer is composed into the policy, the cases it
    distinguishes (missing cache tiers, unavailable load, untracked prefill) are handled exactly as
    the default handles them, and the host's own score-weight settings do not change feature 0.
    % src: FEATURES.md row 0; tests.rs feature_zero_ignores_the_hosts_score_weights (200 random tables)
  \item \textbf{Cache, load and session.} Features 1--5 come from the router's cache estimate and
    its own request tracking; feature 6 needs a bounded session map, because replay never sets the
    router's affinity target (\cref{app:policy}).
    % src: FEATURES.md "Feature set v1" preamble, rows 1, 4, 6; session_map.rs
  \item \textbf{Prompt length.} A request quantity alone, such as $\ISL_j$, is the same for every
    candidate and cannot change an argmax. Prompt length matters only through features that
    combine it with a worker quantity (features 0, 1, 2 and 7; feature 7 is the explicit product
    of prompt length and queued prefill) or through the context term (\cref{sec:context}).
\end{itemize}

% src: CR/facts/pilot.json feature_scaling.pooled_within_set_sd (default_logit_scaled 6.6074,
%      active_requests_s 0.060309; caveat: tables reconstructed approximately from default@defaults
%      per-request rows on the pilot's train subset)
The features' within-set spreads differ by two orders of magnitude, from 0.060 to 6.6 on candidate
tables reconstructed from the default router's replays in the pilot, a smaller tuning run that
preceded the main one (\cref{sec:pilot}; \emph{preliminary (pilot, train subset)}). Training
therefore bounds each coefficient relative to its feature's spread (\cref{sec:space}).

\paragraph{Reading the coefficients.}
% Derivation from eq. (default-cost) and tab:features; no campaign number involved.
With $\theta_0 = -1$, as in every learned arm of the model ladder (M1, M2 and their variants;
\cref{sec:ladder}),
$\util_i = -\cdef_i/\reqblocks_j + \sum_{f \ge 1}\theta_f x_{i,f}$, so every term is measured
in default-cost blocks per request block. Two consequences help when reading tuned values. When
the clip in \cref{eq:default-cost} is inactive, that is, when the overlap credit does not exceed
the prompt's blocks plus the queued prefill blocks, $\cdef_i$ contains $-\ovl_i$, so
$\theta_1 x_{i,1} = \theta_1\ovl_i/\reqblocks_j$ adds $\theta_1$ blocks of credit per
overlapping block to the default's one: $\theta_1 > 0$ is
more cache-seeking than the default. And features 0 and 1 are per request block while features
2--7 are not, so at fixed $\theta$ the anchor's differences between candidates shrink as
$1/\reqblocks_j$ relative to the load features: the balance between the default cost and the
learned load terms depends on prompt length by construction.

\paragraph{Feature set \code{v2}.}
% src: FEATURES.md "Feature set v2"; features.rs (LOG_TOKEN_SCALE 1024, SET_RELATIVE floors,
%      HASH_HOME_PREFIX_TOKENS 256); CR/CONTRACT.md A17.1 (M1-v2 mandatory tier B arm)
\Cref{tab:features-v2} lists 15 features that \code{v2} appends to the eight of \code{v1}, which
keep their indices. They are used by the M1-v2 arm (\cref{sec:ladder}). They add three things
\code{v1} cannot express. Log-scale load terms make products of loads linear in the utility, so
LMetric's multiplicative rule \citep{zhang2026lmetric} lies inside the model class:
$\theta = -(\basis{8} + \basis{10})$ routes to the minimum of
$(\Lpf_i + (\ISL_j - \hit_i)^+)(\Lreq_i + 1)$, LMetric's product with queued prefill, and
$\theta = -(\basis{9} + \basis{10})$ reproduces the branch's \code{lmetric} port as evaluated
(\cref{sec:ports}) without its hot-spot filter when cached tokens are whole overlapping blocks, as in
replay.\footnote{A unit test checks both representations on \num{1000} random tables.}
% src: tests.rs v2_represents_the_lmetric_products (1,000 trials, > 1,500 tie-free checks);
%      FEATURES.md "Representability"; CR/literature/LESSONS.md LR-07
A prefix-hash home indicator marks two hash-chosen candidates per
prompt prefix, the consistent-hash candidate membership of DualMap and Lodestar
\citep{yuan2026dualmap,lim2026lodestar}; \cref{app:policy} discusses its key.
Set-relative load terms make a worker's load count relative to the other candidates; they enter
nonlinearly, because a linear set-relative term would be a no-op (\cref{sec:context}). The ratio
form follows divisive normalization \citep{webb2021normalization}. The product of new prefill and
active requests stands for the prefill--decode interference that two routers model
\citep{jain2024intelligent,srivatsa2025preble}, and the attention term is SMetric's prefill cost
\citep{wang2026smetric}.
% Inference (ours): M1-v2 pins theta_0 = -1, so LMetric's rule is approached as the log-term
% coefficients grow relative to the anchor, not reached exactly; CONTRACT A17.1 s3 init is
% "LMetric's representable point ... scaled to the pinned anchor".

\begin{table}[tbp]
  \centering\small
  \caption{Features that \code{v2} appends to \code{v1}. $(y)^+ = \max(y, 0)$;
    $(\fbar)_f$ is the candidate mean of feature $f$. Indices 12--20 apply one transform to each of
    features 4, 3 and 5, in that order; the ratio floors are 0.01, $512/\Tch$ and $1/32$, about one
    unit of each load.}
  \label{tab:features-v2}
  \begin{tabularx}{\linewidth}{@{}rlY@{}}
    \toprule
    $f$ & Name & $x_{i,f}$ \\
    \midrule
    8 & \code{log\_ptok} & $\ln\!\big(\max(\Lpf_i + (\ISL_j - \hit_i)^+,\,1)/1024\big)$ \\
    9 & \code{log\_new\_prefill} & $\ln\!\big(\max((\ISL_j - \hit_i)^+,\,1)/1024\big)$ \\
    10 & \code{log\_bs} & $\ln(1 + \Lreq_i)$ \\
    11 & \code{hash\_home} & 1 if $i$ has one of the two highest rendezvous weights for the
      request's prefix key (the hash of the prompt's first 256 tokens); 1 for all candidates when
      $|\cand| \le 2$; 0 for all without prefix hashes \\
    12--14 & \code{*\_ratio} & $x_{i,f} / \big(\text{floor}_f + (\fbar)_f\big)$ \\
    15--17 & \code{*\_below\_frac} & $|\{i' \in \cand : x_{i',f} < x_{i,f}\}| / |\cand|$ \\
    18--20 & \code{*\_excess} & $\big(x_{i,f} - (\fbar)_f\big)^+$ \\
    21 & \code{new\_prefill\_x\_requests} & $x_{i,2}\, x_{i,5}$ \\
    22 & \code{prefill\_attn} & $x_{i,2}\,(\ISL_j/\Tch - x_{i,2}/2)$, the attention work of
      prefilling the uncached suffix, in units of $\Tch^2$ \\
    \bottomrule
  \end{tabularx}
\end{table}
% src: features.rs compute() (log features, PREFILL_ATTN = n (L - n/2) / T^2 with n = T x_2, L = ISL),
%      set_relative() (count = |S|, strict below), hash_home() (threshold = second-highest weight;
%      <= 2 candidates -> all home), hash_home_key() (prefix hash at 256 / block_size blocks)

% prov: A16 (v3 feature set). The v3 definitions and their src comments moved to
% sections/appendix-ais-league.tex (app:ais-league).
A third feature set, \code{v3}, adds \ais{} estimates for the secondary \ais{}-informed league,
which never enters the headline (\cref{sec:res-ais}).

\subsection{Utility and decision rule}
\label{sec:utility}

% src: FEATURES.md "Utility and decision"; CR/CONTRACT.md "Utility of candidate i in set S";
%      learned_choice/mod.rs (effective_weights, decide)
For request $j$ with candidates $\cand$ and features $\fvec_i$, the policy scores
\begin{equation}
  \label{eq:utility}
  \util_i = \theta^{\top}\fvec_i + \sum_{\cidx=1}^{\rank} \ctxv_{\cand,\cidx}\, p_{\cidx}^{\top}\fvec_i
          = \weff^{\top}\fvec_i,
  \qquad \weff = \theta + \sum_{\cidx=1}^{\rank} \ctxv_{\cand,\cidx}\, p_{\cidx},
\end{equation}
and decides
\begin{equation}
  \label{eq:decision}
  \temp = 0:\ \ i^{\star} \in \argmax_{i \in \cand} \util_i \ \text{(seeded tie-break)};
  \qquad
  \temp > 0:\ \ \Pr[i^{\star} = i] = \frac{\exp(\util_i/\temp)}{\sum_{i' \in \cand}\exp(\util_{i'}/\temp)} .
\end{equation}
With $\rank = 0$ this is a conditional logit, the simplest member of the generalized
extreme-value family of choice models \citep{train2009gev}: one coefficient vector shared by every
candidate, so the same parameters apply at any $\Nw$. The
context values $\ctxv_{\cand,\cidx}$ are scalars computed once per decision
(\cref{sec:context}); folding them into $\weff$ is what keeps a decision at $O(\Nw\dfeat)$. A
non-finite utility fails the selection rather than routing on it.

M1, M2 and their variants are trained and evaluated at $\temp = 0$, and a $\temp > 0$ variant is
only a conditional robustness extra.\footnote{The temperature is not range-normalized, unlike the
  default's \code{router\_temperature}, which divides costs by their range before scaling; at a
  fixed $\temp > 0$ the probability mass on workers other than the best grows with $\Nw$, so a
  sampled policy does not mean the same thing at every worker count.}
% src: FEATURES.md "Utility and decision" (temperature not range-normalized, LR-15);
%      default/picker.rs softmax_sample_index; CR/literature/LESSONS.md LR-05 (flat directions), LR-15 action 3
At $\temp = 0$, scaling $\theta$ and the $p_\cidx$ together never changes a decision, which is why
training pins $\theta_0 = -1$ (\cref{sec:space}).
% src: CR/facts/gate.json full_plan.tiers.C_conditional ("robustness extras (speedup 0.8/1.2, tau > 0,
%      hash_home)"); CR/literature/LESSONS.md LR-15 action 3
\Cref{alg:decide} in \cref{app:policy} gives one decision in full.

\paragraph{Determinism and tie-breaking.}
% src: FEATURES.md "Determinism"; choice.rs module docs; CR/literature/LESSONS.md LR-03;
%      tests.rs one_draw_per_decision_keeps_later_draws_aligned
Candidates are visited in worker-ID order, and under the default \code{one\_draw} discipline every
decision consumes exactly one random draw, so a seeded policy decides identically across processes
and common random numbers stay aligned across the CMA-ES candidates of a generation
(\cref{app:policy}).

\subsection{Context terms and the IIA restriction}
\label{sec:context}

% src: WT/notes/learned-routing/campaign/PLAN.md M2 row ("captures interactions between
%      workers at any N"); CR/literature/notes/dcm-context.md §1 (LCL), §2.1 (Lemma 4.3);
%      CR/literature/LESSONS.md LR-05, LR-06
With $\rank = 0$, the policy inherits the conditional logit's independence of irrelevant
alternatives (IIA). At $\temp > 0$,
\begin{equation}
  \label{eq:iia}
  \frac{\Pr[i^{\star} = i]}{\Pr[i^{\star} = i']} = \exp\!\big(\theta^{\top}(\fvec_i - \fvec_{i'})/\temp\big),
\end{equation}
which depends only on the two workers' own features, and at $\temp = 0$ the same holds for which of
the two ranks higher. The exchange rate between features, for example how many queued prefill
tokens one cached block is worth ($\theta_1$ against $\theta_3$), is then the same whether the
other workers are idle or saturated. A router plausibly wants it to change: when every candidate
is busy, waiting behind the worker that holds the prefix may beat recomputing the prefix
elsewhere, and when idle workers exist it may not. This is a hypothesis about the best policy,
not an established property; the drift check described below measures it.

A set statistic added directly to the utility cannot express this. A term that is the same for
every candidate, or a linear set-relative feature such as $x_{i,f} - (\fbar)_f$, shifts every
utility in $\cand$ by the same amount and leaves the argmax unchanged. In the linear context
logit, choice likewise depends only on each item's deviation from the set mean
\citep[Lemma~4.3]{tomlinson2021lcl}. Set information has to change the coefficients instead, and
the context term in \cref{eq:utility} does exactly that: $\weff = \theta + \sum_\cidx
\ctxv_{\cand,\cidx}\,p_\cidx$, so the exchange rate between two features becomes a linear function
of the context values. This is the linear context logit (LCL) of \citet{tomlinson2021lcl},
$\util_i = (\theta + A\,\ctxv_{\cand})^{\top}\fvec_i$, with two changes: the mean feature vector is
replaced by a short vector of named set and request statistics (\cref{tab:sources} in
\cref{app:policy}), and $A$ is restricted to one free column $p_\cidx$ per statistic. It is not
their decomposed variant (DLCL), which is a mixture of logits.

\paragraph{When context should pay.}
% src: CR/literature/LESSONS.md LR-15 (Tomlinson footnote 3; drift check); CR/facts/STATE.md
%      phase2-prep ("4 LR-15 drift runs at 96 evals"); CR/facts/gate.json full_plan.tiers;
%      CR/facts/tierA.json drift (final, with its interpretation)
Context helps only where the best coefficients drift with the set statistics, and a larger nested
model can lose on a fixed budget \citep{tomlinson2021lcl}. A drift check before M2 was
inconclusive (\cref{sec:drift}), so it did not choose M2's sources; they were fixed before any M2
evaluation: \code{isl\_k}, because prompt length is M1's documented failure axis
(\cref{sec:gaming}), and \code{mean\_active\_prefill\_tokens\_k}, the candidate set's prefill
pressure.
% The check's numbers (cross-play loss 0.087 on 3 cells, 1.3 SE; a single model beat every regime
% specialist) are in sec:drift (appendix-tuning.tex); restructure round 1, E2.
% src: CR/facts/tierA.json drift (cross_play_losses: l1-tuned on l3 0.0865; interpretation: verdict
%      INCONCLUSIVE, evidence_of_drift false, use_for_m2_context_sources false, why[0..3]: SE 0.0645 over
%      3 cells, distance moved 0.08-0.28 vs 0.83, pooled beats every specialist);
%      CR/facts/tierBC.json decisions.m2_rank2 (sources, why_sources)

\subsection{Set-size invariance}
\label{sec:set-size}

% src: CR/cells/SPLIT_MANIFEST.json (train N {4,8}; test adds 2, 16, 32); FEATURES.md "Context sources"
%      (duplication paragraph); CR/literature/LESSONS.md LR-06; tests.rs
%      duplicating_the_worker_set_keeps_features_sources_and_choice (50 random tables, v2, random theta,
%      theta[hash_home] = 0, requests without a session)
Training sees $\Nw \in \{4, 8\}$; the test also scores $\Nw = 2$, 16 and 32 (\cref{sec:splits}).
Two properties make one coefficient vector meaningful at any worker count. First, the
coefficients are shared by all candidates and no feature is a worker's identity, so the policy is
defined for any number of candidates; Lodestar's router uses the same principle
\citep{lim2026lodestar}. Second, every feature and context source is designed to be invariant
under duplicating the candidate set: if every worker is replaced by two identical copies
($\Nw \to 2\Nw$), each candidate's features, every context value and the chosen worker (up to
which copy) stay the same. Means and maxima pass this test; a mean of a one-hot feature, which is
$1/\Nw$, does not. A unit test duplicates random \code{v2} candidate tables and checks that
features, sources and the chosen worker are unchanged (\cref{app:policy}). Only \code{hash\_home}
and \code{session\_affinity} are exempt, because they are relations to particular workers.
% src (LR-12): CR/literature/LESSONS.md LR-12 ("Matched per-worker load does not hold the routing
%      regime fixed across N")

Duplication invariance is necessary for transfer across worker counts, not sufficient. At a
different $\Nw$, the same load level produces different per-worker states, so the held-out worker
counts in the test measure transfer directly (\cref{sec:eval}).

\subsection{Parity with the default router}
\label{sec:parity}

% src: FEATURES.md "Parity anchor"; tests.rs theta0_picks_the_default_argmin (2,000 tables; shapes:
%      full, no tier overlaps, no worker loads, prefill untracked; >= 1,500 tie-free; v1 and v2 via YAML)
With $\theta = -\basis{0}$, $\rank = 0$ and $\temp = 0$, \cref{eq:utility,eq:decision} pick
$\argmax_i (-\cdef_i/\reqblocks_j) = \argmin_i \cdef_i$, the default router's choice. Unit tests
check through the deployment YAML path that \lc{} at $-\basis{0}$ picks the default's unique
minimum on 2,000 random tables (\cref{app:policy}).
% src: CR/facts/integration.json parity (11 cells x 8 replicates = 88 pairs; reservoir: identical
%      per_request sha256 88/88; one_draw: identical rows 0/88; noise rule not exceeded)
In replay, the \code{reservoir} tie-break, which replays the seeded default picker's random draws,
reproduces the default's per-request rows on 88 of 88 (cell, replicate) pairs; the
\code{one\_draw} tie-break used in training breaks ties differently but changes goodput by no more
than replicate noise (mainly on Mooncake cells; the session and AgentX cells sat at ceiling under the
provisional SLO used at integration; \cref{app:policy}).
% src (qualifier): CR/facts/integration.json parity caveat (sessions and AgentX cells near good
%      fraction 1.0 at the provisional SLA; 7 Mooncake cells carry the test); restructure round 1, F5

This parity is what makes $-\basis{0}$ a meaningful start for training (restart s1 in
\cref{sec:restarts}) and lets every learned coefficient be read as a deviation from the default
router.

\subsection{Model ladder}
\label{sec:ladder}

% src: WT/notes/learned-routing/campaign/PLAN.md "Learned model ladder"; CR/facts/gate.json
%      full_plan.tiers; CR/CONTRACT.md A15, A16, A17; CR/facts/pilot.json decisions (M0 space: 5 knobs)
Each rung adds capacity only where the previous one leaves a measured gap, and a rung counts only
if it improves held-out results. \Cref{tab:ladder} lists the arms that were trained. Every arm
received three restarts of the same budget (\cref{sec:budget}). Training and initialization are in
\cref{sec:training}; the rule that picks the headline learned model among the router-observable
arms, the same rule that picks the best baseline, is in \cref{sec:selection}.
% prov: tiers A/B/C. The main tuning scheduled the arms in tiers: tier A held M1 and every tuned
%      baseline, tier B the arms that start from the selected M1 or were added by later amendments,
%      and tier C the arms that ran only if an earlier result called for them.
% src: CR/facts/gate.json full_plan.tiers (A_headline_core, B_ladder, C_conditional);
%      CR/facts/tierA.json, CR/facts/tierBC.json (status complete; 36 + 30 runs)

\begin{table}[tbp]
  \centering\small
  \caption{The learned arms and the tuned default. ``Free'' counts searched coefficients. Every
    learned arm runs at $\temp = 0$ with $\theta_0$ pinned at $-1$ and, after the gaming finding,
    with load coefficients clamped so that load cannot attract (\cref{sec:sign-constraint}). Only
    router-observable arms can be the headline learned model; results are in
    \cref{sec:res-ladder}.}
  \label{tab:ladder}
  \begin{tabularx}{\linewidth}{@{}llrY@{}}
    \toprule
    Arm & Features & Free & Role and outcome \\
    \midrule
    M0 & \defaultcf{} & 5 & tuned default: $\bov$, $\bdecay$, $\bpf$, $\breq$ and the router
      temperature \\
    M1 & \code{v1}, $\rank = 0$ & 7 & headline candidate; sign-constrained re-run \\
    % prov: "sign-constrained re-run of tier A"
    M1 unconstrained & \code{v1}, $\rank = 0$ & 7 & the first tuned M1; flagged for gaming,
      reported as a secondary row \\
    % prov: "the original tier-A M1; gaming-flagged"
    M1 continued & \code{v1}, $\rank = 0$ & 7 & 400 more evaluations from the selected M1; the
      comparator for M2 \\
    M2 rank 2 & \code{v1}, two named sources & 20 & headline candidate; sources \code{isl\_k} and
      \code{mean\_active\_prefill\_tokens\_k}; ranks 3--4 not run (their pre-registered gate
      failed) \\
    M1-noaff & \code{v1}, $\theta_6 = 0$ & 6 & headline candidate; session coefficient pinned
      at 0 \\
    % prov: A15
    M1 + queue threshold & \code{v1} and \code{router\_queue\_threshold} & 8 & headline candidate
      (queue-threshold ablation) \\
    % prov: ablation (a)
    M1-v2 & \code{v2}, $\rank = 0$ & 22 & headline candidate; extended features; selected \\
    % prov: A17.1
    M3 & nested logit & -- & not run: no nontrivial nests (\cref{sec:nested}) \\
    \midrule
    M1-ais, M2-ais & \code{v3} (and two named sources) & 11, 33 & secondary \ais{}-informed
      league \\
    % prov: A16
    \bottomrule
  \end{tabularx}
\end{table}
% src: free counts from CR/runs/phase2/spaces/{m0_default_tuned,m1c_learned_v1,m1_learned_v1,
%      m1_continue,m2_rank2,m1noaff_s1,abla_m1,m1v2_s1}.yaml and CR/runs/phase2-ais/spaces/
%      {m1c_ais_s1,m2_ais_rank2}.yaml (bounds with low != high); roles: CR/facts/finalists.json
%      (headline candidates m1, m2r2, m1noaff, m1v2, ablam1; m1cont and m1u secondary);
%      CR/facts/tierBC.json decisions.m2_rank2 (sources), m2_rank3_4_gate

The later router-observable arms test specific questions. M1-noaff pins the session coefficient at
zero: the hypothesis was that session affinity carries little beyond cache overlap in
replay, where the router's cache index is perfectly fresh, but may matter under router-state lag
or live, so the arm also ran in the lag robustness check. M1-v2 targets the pilot's weakest case,
the agent workloads on which M1 trailed cache-aware heuristics, with \code{v2}'s log, set-relative
and prefill-attention terms; one of its starts is LMetric's representable point. Because $\theta_0$
stays pinned at $-1$, that point is approached as the log-term coefficients grow relative to the
anchor rather than reached exactly (our inference from \cref{tab:features-v2}).
% prov: "the operator's hypothesis" (A15); A17.1.
% src: CR/CONTRACT.md A15 ("Hypothesis: affinity is near-redundant in replay, where the index is
%      perfectly fresh, but may matter under router-state lag or live"); A17.1 ("Why: ... M1 trailing on
%      AgentX. v2's log, set-relative and prefill-attention terms target exactly that")

\label{sec:ablations}
Two ablations were planned, and neither ships. The queue-threshold ablation tuned the router's
dispatch-or-hold threshold \code{router\_queue\_threshold} jointly with the selected learned model
and with the best baseline; it ran, and found nothing measurable
(\cref{sec:res-ladder}).\footnote{The other ablation, adding simulator-only signals such as engine
  waiting counts and KV usage, was not run because of its implementation cost (\cref{app:tuning}); the
  \ais{}-informed league measures part of the same question (\cref{sec:res-ais}).}
% prov: ablation (a) = the queue-threshold ablation; ablation (b) = simulator-only signals.
% src: CR/facts/tierBC.json decisions.ablation_a, tier_c_final_i9, ablation_b, ablation_b_recheck_i8;
%      CR/CONTRACT.md A17.3 (cut order: ablation b first); CR/report/REPORT.md section 5
